\chapterbackmatter{Solutions to Weinberg Exercises}
\ExerciseEditorialNote

\begin{enumerate}
  \WeinbergSolution{1}
  Let \(\phi_i\) be real elementary scalars in a (possibly reducible)
  representation of the gauge algebra.  It is useful to choose
  antisymmetric Hermitian generators \(T_a\), including the coupling
  in the same way as the chapter's \(t_a\), and to add to the BRST
  transformations
  \[
    s\phi_i=i\omega^a(T_a)_i{}^j\phi_j,\qquad
    [T_a,T_b]
      =iC_{cab}T_c .
  \]
  The Jacobi identity and the representation relation give
  \(s^2\phi_i=0\).  Introduce a source \(K_{\phi i}\) for \(s\phi_i\).
  Since \([\phi]=1\), it has
  \[
    [K_{\phi}]=2,\qquad
    {\rm gh}(K_{\phi})=-1,
  \]
  and statistics opposite to those of \(\phi\).

  The loopwise Zinn--Justin argument is unchanged.  Once subdivergences
  through \(N-1\) loops have been subtracted, locality and power counting
  imply that \(\Gamma_{N,\infty}\) is an integrated polynomial of
  dimension at most four.  Its antifield-dependent part is at most linear
  in \(K_\phi\): a term quadratic in \(K_\phi\) already has dimension four
  and ghost number \(-2\), while no field of positive ghost number can be
  added without exceeding dimension four.  Mixed quadratic terms are
  excluded by the same dimension and ghost-number counting used in
  Section 17.2.  Thus
  \[
    \Gamma_{N,\infty}
      =\Gamma_{N,\infty}[\chi,0]
       +\int d^4x\,K_n\mathcal D_N^n ,
    \qquad
    (S_R,\Gamma_{N,\infty})=0 .
  \]
  The last equation says that the action and the BRST transformation are
  infinitesimally deformed together, and that the deformed transformation
  is nilpotent.

  Lorentz covariance, ghost number, and dimension now restrict the scalar
  part of that deformation to
  \[
    s'\phi_i=i\mathcal Z\omega^a
       (T'_a)_i{}^j\phi_j ,
  \]
  apart from a possible transformation of a constant scalar shift when
  the fields are expanded about a vacuum.  Nilpotence requires
  \[
    [T'_a,T'_b]
       =iC'_{cab}T'_c .
  \]
  Hence \(T'_a\) is the same representation after a change of scalar
  basis and the common gauge-coupling renormalization already found for
  the gauge and fermion sectors.  For repeated equivalent irreducible
  representations, the allowed change of basis is a matrix wave-function
  renormalization among the copies.

  It remains to classify the scalar-dependent part of the local action.
  Before gauge fixing, the most general renormalizable expression is
  \[
  \begin{split}
    \mathcal L_{\phi}={}&
      -\frac12 Z_{ij}(D_\mu\phi)_i(D^\mu\phi)_j
      -\ell_i\phi_i-\frac12m^2_{ij}\phi_i\phi_j\\
      &-\frac1{3!}\kappa_{ijk}\phi_i\phi_j\phi_k
      -\frac1{4!}\lambda_{ijkl}\phi_i\phi_j\phi_k\phi_l ,
  \end{split}
  \]
  supplemented, when fermions are present, by every gauge-invariant
  Yukawa tensor \(Y^i\).  Gauge invariance says, for example,
  \[
    (T_a)_i{}^{i'}\lambda_{i'jkl}
    +(T_a)_j{}^{j'}\lambda_{ij'kl}
    +(T_a)_k{}^{k'}\lambda_{ijk'l}
    +(T_a)_l{}^{l'}\lambda_{ijkl'}=0,
  \]
  with analogous equations for \(\ell,m^2,\kappa\), and \(Y\).  These are
  precisely the tensors available as counterterms.

  Finally, scalar fields enlarge the renormalizable gauge-fixing fermion.
  Depending on the gauge, it may contain
  \(\omega_a^*u_{a i}\phi_i\) or
  \(\omega_a^*u_{a ij}\phi_i\phi_j\); after applying \(s\) these
  generate the \(h\phi\) and \(h\phi\phi\) terms mentioned in
  Section 17.2, together with their required ghost partners.  They are
  BRST-exact and are absorbed by renormalizing the gauge-fixing fermion.
  Therefore all possible divergences are removed by scalar
  wave-function, potential, Yukawa, gauge-coupling, and BRST-exact
  gauge-fixing counterterms.  The induction closes, provided the matter
  representation is anomaly free and the regulator preserves, or is
  corrected to restore, the Zinn--Justin identity.

  \WeinbergSolution{2}
  Write the full gauge field as
  \(\mathcal A_{a\mu}=A_{a\mu}+A'_{a\mu}\), where \(A\)
  is the prescribed background and \(A'\) is integrated over.  The
  quantum BRST differential leaves the background fixed and is obtained
  from the true gauge transformation \eqref{eq:17.4.17} by replacing its
  parameter by the ghost:
  \[
  \begin{aligned}
    sA_{a\mu}&=0,\\
    sA'_{a\mu}
      &=\bar D_\mu\omega_a
        -C_{abc}\omega_b A'_{c\mu},\\
    s\omega_a
      &=-\frac12C_{abc}\omega_b\omega_c,\\
    s\omega_a^*&=-h_a,\qquad &sh_a&=0 .
  \end{aligned}
  \]
  Here \(\bar D_\mu\) contains only \(A\).  On a quantum matter
  fluctuation,
  \(s\psi'=it_a\omega_a(\psi+\psi')\), while
  \(s\psi=0\).  These rules are nilpotent because their action on the
  sums \(A+A'\) and \(\psi+\psi'\) is the ordinary Yang--Mills BRST
  differential.

  Choose the background-covariant gauge-fixing fermion
  \[
    \Psi=-\int d^4x\,\omega_a^*
       \left(\bar D^\mu A'_{a\mu}+\frac{\xi}{2}h_a\right).
  \]
  Up to the harmless overall sign convention for \(s\), its variation is
  \[
  \begin{split}
    s\Psi=\int d^4x\Bigl\{&
       h_a\bar D^\mu A'_{a\mu}
       +\frac{\xi}{2}h_a h_a\\
       &+\omega_a^*\bar D^\mu
       [\,\bar D_\mu\omega_a
          -C_{abc}\omega_b A'_{c\mu}\,]
       \Bigr\}.
  \end{split}
  \]
  Eliminating \(h_a\) gives
  \(-(\bar D^\mu A'_{a\mu})^2/(2\xi)\), while the last line is
  precisely the background-gauge Faddeev--Popov operator.  Thus the
  gauge-fixed action is
  \[
    S_{\rm gf}=I_0[A+A',\psi+\psi']+s\Psi .
  \]
  Adding sources \(K_n\) for the nonlinear quantum BRST variations gives
  the usual Zinn--Justin identity \((\Gamma,\Gamma)=0\), so the direct
  renormalization argument applies even though antighost translation is
  no longer available; one simply admits the full set of local
  BRST-exact counterterms \(s\Psi_{\rm ct}\).

  There is in addition a background gauge transformation.  The
  background \(A\) transforms as a connection, while
  \(A'\), \(\omega\), \(\omega^*\), and \(h\) transform homogeneously in
  the adjoint representation (and both background and quantum matter
  transform in their matter representation).  Consequently
  \(\bar D^\mu A'_\mu\) transforms homogeneously, \(\Psi\) and \(s\Psi\)
  are invariant, and the functional measure may be chosen invariant.
  The effective action of the background fields therefore obeys an
  ordinary background Ward identity.

  In particular, its divergent pure-background term has the form
  \[
    \Gamma_\infty[A]\supset
      -\frac14L_A\int d^4x\,
       F_{a\mu\nu}[A]F_a^{\mu\nu}[A].
  \]
  Defining \(A^R=(1+L_A)^{1/2}A\), background gauge invariance requires
  the product of the coupling and the background connection to remain
  fixed:
  \[
    g^R=g(1+L_A)^{-1/2},
    \qquad Z_gZ_A^{1/2}=1 .
  \]
  This is the practical advantage of the construction: the coupling
  counterterm is obtained from the background two-point function alone.

  \WeinbergSolution{3}
  Let the background field carry momentum.  Background gauge invariance
  makes its one-loop two-point vertex transverse:
  \[
    \Gamma^{(1)}_{AA}
      =\frac12\int\frac{d^dq}{(2\pi)^d}\,
       A_\mu^a(-q)
       (q^2\eta^{\mu\nu}-q^\mu q^\nu)
       \Pi(q^2)A_\nu^a(q).
  \]
  In background Feynman gauge the diagrams are the quantum-gauge loop
  (including its seagull), the ghost loop, and the matter loop.  Tensor
  reduction may be performed after combining denominators with
  \[
    \frac1{k^2(k+q)^2}
      =\int_0^1dx\,
       \frac1{\{(k+xq)^2+x(1-x)q^2\}^2}.
  \]
  Terms proportional to \(\eta_{\mu\nu}\) that are not transverse cancel
  between the bubble and seagull graphs.  The remaining pole and
  logarithm have the group factor
  \[
    b_0=\frac{11}{3}C_1-\frac{4}{3}C_2 ,
  \]
  where the definitions of \(C_1\) and \(C_2\) are those of
  Eqs.~\eqref{eq:17.5.33} and \eqref{eq:17.5.34}.  Equivalently, the local
  part of the answer is
  \[
    \Gamma^{(1)}_{AA,\infty}
      =-\frac14L_A\int d^4x\,F_{a\mu\nu}F_a^{\mu\nu},
  \]
  with
  \[
    L_A=
      \frac{g^2}{2\pi^2}
      \left(\frac{11}{12}C_1-\frac13C_2\right)
      \left[\frac1{d-4}+\ln\mu+\cdots\right]+O(g^4).
  \]
  The same coefficient multiplies
  \(-\ln(-q^2)\) in the nonlocal part; changing the infrared reference
  scale only shifts the displayed local logarithm.  The gauge and ghost
  fields together supply \(11C_1/3\), whereas a Dirac matter multiplet
  supplies \(-4C_2/3\).  This separation is gauge dependent, but their
  sum is not.

  Finally use the background Ward-identity relation
  \(g_R=g(1+L_A)^{-1/2}\).  Expanding the square root gives
  \[
    \boxed{
    g_R=g\left[1-\frac{g^2}{4\pi^2}
      \left(\frac{11}{12}C_1-\frac13C_2\right)
      \left(\frac1{d-4}+\ln\mu+\cdots\right)
      +O(g^4)\right],}
  \]
  which is Eq.~\eqref{eq:17.5.44}.  Thus the spacetime-dependent
  two-point calculation and the constant-background four-point
  calculation determine the same gauge-invariant operator \(F^2\).

  \WeinbergSolution{4}
  Let the scalar generators, including the coupling \(g\), be normalized
  for each complex scalar multiplet \(r\) by
  \[
    {\rm tr}_r(T_a T_b)
       =g^2 C_s(r)\delta_{ab}.
  \]
  There are two one-loop diagrams quadratic in the background: the
  scalar bubble with two derivative vertices and the seagull with one
  \(A^2\phi^*\phi\) vertex.  Their non-transverse parts cancel.  For one
  complex multiplet the transverse pole has one quarter of the magnitude
  of a Dirac multiplet of the same index, and hence contributes
  \(-C_s(r)/3\) to \(b_0\).  A real multiplet has half as many independent
  modes and contributes \(-C_s(r)/6\).

  For arbitrary Dirac fermions, complex scalars, and real scalars,
  therefore,
  \[
    b_0=\frac{11}{3}C_1-\frac{4}{3}C_2
       -\frac13\sum_{\text{complex }r}C_s(r)
       -\frac16\sum_{\text{real }r}C_s(r).
  \]
  The chapter's bracket is simply \(b_0/4\).  Thus dimensional
  regularization gives
  \[
    L_A=\frac{g^2b_0}{8\pi^2}
      \left[\frac1{d-4}+\ln\mu+\cdots\right]+O(g^4)
  \]
  and
  \[
    \boxed{
    g_R=g\left[1-\frac{g^2b_0}{16\pi^2}
      \left(\frac1{d-4}+\ln\mu+\cdots\right)+O(g^4)\right].}
  \]
  With a hard ultraviolet cutoff the same result reads
  \[
    g_R=g\left[1+\frac{g^2b_0}{16\pi^2}
      \ln\frac{\Lambda}{\mu}+O(g^4)\right].
  \]
  As a check, for \(SU(N)\) with \(n_f\) Dirac fundamentals and \(n_s\)
  complex scalar fundamentals, \(C_1=N\) and
  \(C_2=n_f/2\), \(C_s=1/2\), so
  \[
    b_0=\frac{11}{3}N-\frac23n_f-\frac16n_s .
  \]
\end{enumerate}
